% Transcription — fonds Alexandre Grothendieck, shelfmark 98, batch 2 (pages 21-40).
% Established from the facsimile archives/batches/98/batch-02.pdf (a mirror of
% https://grothendieck.umontpellier.fr/98.pdf), per the skill
% transcribe-grothendieck, read from the tiles under archives/tiles/98/ and
% from 300-600 dpi crops of the batch PDF. The batch file carries no cover
% sheet: PDF page k is folder page 20+k, and the archivists' pencil numbers
% (« 23 », « 26 » ... « 40 ») agree wherever legible.
%
% Pass: Opus 5.5 (claude-opus-5-5), 2026-09-24 — first pass, unchecked against the pages by a human.
%
% The folder sits in the inventory group "69-89" « Géométrie et topologie
% combinatoire », whose subject does not match this folder's (as for folder
% 100); \dating{} carries the folder's own date verbatim. Nothing is
% concluded from the group.
%
% Seventeen of the twenty pages are transcribed: 21, 23, 26-40. Absent:
%   - pages 22 and 24: typed leaves of someone else's typescript (French,
%     scanned upside down, headed « n° 268 », paginated « -87- » and
%     « -95- »), on abelian and Picard varieties. Neither carries a mark in
%     his hand and neither bears on these notes: skipped under the rule of
%     folders 100 and 71 (other people's text with no mark of his);
%   - page 25: an otherwise blank sheet carrying, in his hand, « Familles
%     limitées » and « cf. SGA 6 XIII ». It is a cover: no \page{}, and his
%     words become the section heading, with a \note.
% No page is summarised; everything transcribed is his manuscript.
%
% What the batch is.
%   21, 23  Two isolated sheets of spectral sequences: the exact sequence of
%           low degree for H^*(-/S, G_m) (Brauer group, d_3), then a table of
%           four fibrations and their spectral sequences (Leray,
%           Leray-Cartan, Hochschild-Serre), H^*(X) <= H^p(Y, H^q(F)) etc.
%   25-27   « Familles limitées »: sheaves F_i quotients of a fixed E with
%           fixed Hilbert polynomials; induction on the dimension of the
%           support via cycles of bounded degree; « Résultat de Chow »
%           (cycles of bounded degree in P^n_Z form a limited family).
%   28-33   k-equivalence of algebraic sheaves defined over different
%           extensions of k (Définition 1, Proposition 1), the set 𝔉(X/k) of
%           classes, its description as a quotient by Aut(K~/k) for K~
%           algebraically closed of infinite transcendence degree; set-valued
%           operations T(F_1,...,F_n) compatible with extension of the base
%           field and the induced 𝔉(A_1,...,A_n); a list of examples
%           (Hom, ⊗, Tor, Ext, extensions, primary components ...).
%   34-40   Stability of limited families: (a) A ⊗ B, Hom(A,B), Tor_i,
%           Ext^i of limited families are limited (proof by stratifying the
%           base), a lemma on a complex of coherent sheaves whose cohomology
%           commutes with base change on a finite stratification, (b) f^*
%           and Tor^{O_Y}(A, O_X); then base change for Tor: a Proposition
%           (dense open U, S' = U_red), a Lemme via the spectral sequence
%           Tor^{B'}_*(M'_1,...,M'_n) <= Tor^A_p(Tor^B_q(M_1,...,M_n), A'),
%           a Corollaire (finite stratification of S), and the start of the
%           same Proposition for Ext, which batch 3 (page 41) continues.
%
% Notation fixed for the batch: underlined Hom, Tor, Ext (sheaf versions)
% are \mathcal{H}om, \mathcal{T}or, \mathcal{E}xt, as in batch 3; his
% gothic F for the set of classes of sheaves is \mathfrak{F}; X~, K~ are
% \widetilde{X}, \widetilde{K}. The hand on pages 26-27 and 31-36 is a fast
% draft: statements and formulas come through, the connective prose is
% often lost and stays \ill{}.
\documentclass[11pt,a4paper]{article}
% Le préambule commun à toutes les transcriptions du fonds Grothendieck.
%
% Il tient en une idée : **l'incertitude se compose, elle ne se résout pas.**
% Une transcription qui lisse un mot douteux a détruit la seule chose pour
% laquelle on la faisait — un lecteur doit pouvoir savoir, à chaque endroit,
% ce qui était sur le feuillet et ce qui a été ajouté. Les six macros qui
% suivent sont donc l'appareil critique, et non de la décoration.
%
% Les mêmes commandes sont reconnues par `scripts/render.mjs`, qui produit la
% vue de lecture du site. Une macro ajoutée ici doit l'être aussi là-bas,
% faute de quoi le rendu échouera bruyamment — ce qui est voulu : un rendu
% silencieusement faux est pire qu'un rendu absent.

\ProvidesPackage{grothendieck}[2026/08/08 Transcriptions du fonds Grothendieck]

\usepackage{amsmath,amssymb,amsthm}
\usepackage{mathrsfs}
\usepackage{stmaryrd}
% Les diagrammes commutatifs sont partout dans ce fonds : c'est la langue
% même de la géométrie algébrique de Grothendieck, et une transcription qui
% les rendrait en prose ne transcrirait rien.
\usepackage{tikz-cd}
% Français par défaut — c'est la langue des feuillets — mais l'anglais est
% chargé aussi : la traduction et le résumé sont des documents anglais, et la
% typographie française (espaces avant les deux-points) y serait une faute.
% Ces documents basculent par \selectlanguage{english} en tête de corps.
\usepackage[english,french]{babel}
\usepackage{fontspec}
\usepackage{geometry}
\geometry{a4paper,margin=25mm,marginparwidth=28mm}
\usepackage{marginnote}
\usepackage{xcolor}
% ulem, not soul. `soul`'s \st and \ul are built on a character-by-character
% scan that XeLaTeX's font machinery breaks: the document fails with an
% undefined control sequence rather than degrading. ulem does the same two jobs
% and is Unicode-engine safe. `normalem` stops it hijacking \emph, which the
% transcriptions use for Grothendieck's own emphasis.
\usepackage[normalem]{ulem}
\usepackage{hyperref}
\hypersetup{colorlinks=true,linkcolor=black,urlcolor=black,pdfborder={0 0 0}}

% --- Métadonnées du lot -----------------------------------------------------
% Reprises telles quelles par le script de rendu, qui les affiche en tête de
% la vue de lecture. Elles fixent ce que le fichier prétend couvrir : sans
% elles, une transcription flotte sans point d'attache dans un fonds de seize
% mille feuillets.

\newcommand{\folder}[1]{\gdef\grothFolder{#1}}
\newcommand{\batch}[1]{\gdef\grothBatch{#1}}
\newcommand{\leaves}[2]{\gdef\grothFirst{#1}\gdef\grothLast{#2}}
\newcommand{\foldertitle}[1]{\gdef\grothFoldertitle{#1}}
\gdef\grothFolder{?}\gdef\grothBatch{?}\gdef\grothFirst{?}\gdef\grothLast{?}\gdef\grothFoldertitle{}

% --- Le résumé --------------------------------------------------------------
%
% Il ouvre la lecture modernisée, sous le titre « Résumé », et il est composé
% en petit. Ce n'est pas un ornement : c'est le seul passage du document qui
% ne soit pas des mathématiques, et un lecteur doit pouvoir le reconnaître
% avant de l'avoir lu — puis le sauter s'il connaît déjà le sujet, ou s'y
% arrêter s'il ne le connaît pas. Le filet qui le suit marque la reprise du
% corps.

\newenvironment{resume}
  {\small\setlength{\parskip}{0.45em}}
  {\par\medskip\hrule\medskip}

% --- Mots-clés ---------------------------------------------------------------
%
% En anglais, en clôture du résumé de la lecture modernisée : le vocabulaire
% moderne sous lequel chercher ce que les feuillets construisent. Le manifeste
% du site (`npm run manifest`) les extrait pour étiqueter la cote entière —
% cette ligne est la source unique des tags, il n'y a pas de fichier à côté.
\newcommand{\keywords}[1]{\par\smallskip\noindent
  {\footnotesize\textsc{Keywords} --- \textit{#1}}\par}

% --- L'appareil critique ----------------------------------------------------

% Le numéro de feuillet, en marge. C'est l'ancre qui permet de régler un
% désaccord sur une formule en regardant *un* feuillet plutôt qu'une chemise
% de six cents. Le site s'en sert aussi pour tourner les pages du fac-similé
% quand on fait défiler la transcription.
\newcommand{\leaf}[1]{%
  \marginnote{\footnotesize\color{gray!70}\textsf{#1}}%
  \label{leaf:#1}\ignorespaces}

% Illisible : on ne devine pas. Un mot inventé qui se lit comme les autres est
% le pire résultat possible.
\newcommand{\ill}{\textcolor{red!60!black}{[\,\ldots\,]}}

% Lecture douteuse : le mot est donné, et signalé comme incertain.
\newcommand{\uncertain}[1]{\uline{#1}}

% Ajout de l'éditeur — une lettre manquante, un mot sous-entendu.
\newcommand{\add}[1]{\textcolor{blue!50!black}{[#1]}}

% Biffé par Grothendieck lui-même. Ce qu'il a rayé fait partie du document :
% une rature dit souvent où la pensée a changé de direction.
\newcommand{\struck}[1]{\sout{#1}}

% Note du transcripteur, sur l'état matériel du feuillet.
\newcommand{\note}[1]{\par\smallskip{\small\color{gray!60!black}\textit{[#1]}}\par\smallskip}

% Note marginale de Grothendieck, distincte du corps du texte.
\newcommand{\marginal}[1]{\par\smallskip\noindent\hspace{1em}%
  {\small\color{gray!60!black}#1}\par\smallskip}

% --- Titre ------------------------------------------------------------------

\AtBeginDocument{%
  \begin{center}
    {\large\bfseries Fonds Alexandre Grothendieck --- cote n\textsuperscript{o}~\grothFolder}\\[2pt]
    {\itshape\grothFoldertitle}\\[4pt]
    {\small feuillets \grothFirst--\grothLast\ (lot \grothBatch)}\\[2pt]
    {\footnotesize\color{gray!60!black}Transcription établie d'après le fac-similé numérisé
     par l'Université de Montpellier.\\
     Les lectures incertaines sont soulignées, les illisibles notées [\,\ldots\,],
     les ajouts entre crochets.}
  \end{center}
  \vspace{1em}}

\endinput


\folder{98}
\batch{2}
\pages{21}{40}
\dating{[à partir de 1966-1967]}
\watermark{Édition de démonstration}
\foldertitle{Schémas de Hilbert. Normes : notes manuscrites (s.d.).}

\begin{document}

\section{[Suites spectrales]}
\note{titre de l'éditeur, entre crochets : les pages 21 et 23 sont deux
feuillets isolés, sans titre, sans lien apparent avec ce qui les entoure.}

\page{21}
$H^{*}(-/S, \mathbf{G}_m) \Leftarrow$
\struck{$\cdot\ H^2(-/S, \mathbf{G}_m)$}

\[
  \cdot \to H^2(T/S, \mathbf{G}_m) \to H^2(-/S, \mathbf{G}_m) \to
  H^0(T/S, \mathcal{H}^2(\mathbf{G}_m)) \xrightarrow{\ d_3\ }
  H^3(T/S, \mathbf{G}_m) \to H^3(-/S, \mathbf{G}_m)'
\]
\[
  \to H^1(T/S, \mathcal{H}^2(\mathbf{G}_m)) \xrightarrow{\ d_3\ }
  H^4(T/S, \mathbf{G}_m) \to E_4^{4,0} \to 0
\]
\note{la première ligne est encadrée ; la seconde est écrite sous le
cadre et la prolonge. Devant $E_4^{4,0}$, une flèche double biffée. Le
signe $\mathbf{G}_m$ de $H^4$ est récrit sur une tache. À gauche, un
schéma $T \leftarrow T \times_S T$ au-dessus de $S$, sans autre texte.}

Au pied de la page, un carré de corps :

\begin{tikzcd}
  k \arrow[r, no head, "\mathrm{rad}"] \arrow[d, no head] & k' \arrow[d, no head] \\
  k'' \arrow[r, no head, "\mathrm{rad}"'] & k'k''
\end{tikzcd}

\note{les quatre côtés sont des traits sans flèche ; les deux côtés
horizontaux portent « rad », les verticaux une marque qu'on lit
\uncertain{sép}. À côté, isolé, $H^2(\ /k, \mathbf{G}_m)$.}

\page{23}
\marginal{\emph{Leray}}

\begin{tikzcd}
  F \arrow[r, hook] & X \arrow[d] \\
  & Y
\end{tikzcd}

application fibrée (au sens de \ill{}) générale
\[
  H^{*}(X) \Leftarrow H^p(Y, H^q(F))
\]

\marginal{\emph{Leray-Cartan}}

\begin{tikzcd}
  X \arrow[d] \\
  Y \arrow[d] \\
  B_F
\end{tikzcd}

$F$ un groupe, fibré principal $\leftarrow$ d'une fibration générale
\[
  H^{*}(Y) \Leftarrow H^p(B_F, \{H^q(X)\})
\]
N.B. $\{H^q(X)\}$ est un système \uncertain{local}
\uncertain{déterminé} par les opérations de $\pi_0(F)$ sur $H^q(X)$.
\note{« \uncertain{déterminé} » est écrit au-dessus d'un mot souligné
qu'on ne lit pas.}

\begin{tikzcd}
  Y \arrow[d] \\
  B_F \arrow[d] \\
  B_{\struck{X}}
\end{tikzcd}

\struck{$X$} $Y$ un groupe, $F$ un s.-groupe, d'une fibration générale
\[
  H^{*}(B_F) \Leftarrow H^p(B_{\struck{X}}, \{H^q(Y)\})
\]
N.B. $\{H^q(Y)\}$ est un système défini par les op.\ de $\pi_0(X)$ sur
$H^q(X/F)$ \ill{} \ill{} \ill{} $\pi_0(X)$ sur $\pi_0(X/F)$.
\note{ici, et dans le schéma au-dessus, il écrit $X$ là où il vient de
nommer le groupe $Y$ ; le $X$ de $B_X$ est biffé, celui du N.B. ne l'est
pas. Laissé tel quel.}

\struck{\ill{}}
\marginal{Sch.\ \ill{}}
Si $Y$ est discret, \ill{} \ill{} d'isomorphismes \ill{} \ill{} de $B_F$,
et telle \ill{} de $B_G$ \ill{} $F$ \ill{} \ill{} \ill{} :
[à \ill{} \ill{} \ill{} $F$, $G$ \ill{} discrets]
\note{trois lignes serrées, d'écriture rapide, sous une ligne biffée ;
seuls les signes mathématiques se lisent. $G$ apparaît ici pour la
première fois.}

\marginal{\emph{Hochschild-Serre}} \marginal{\uncertain{Leray-Čech}}

\begin{tikzcd}
  B_F \arrow[d] \\
  B_{\struck{X}} \arrow[d] \\
  B_Y
\end{tikzcd}

$X$ un groupe, $F$ un sous-groupe distingué, $Y$ le groupe
\uncertain{quotient}, d'une fibration générale
\[
  H^{*}(B_G) \Leftarrow H^p(B_Y, H^q(B_F))
\]
\note{les quatre intitulés sont écrits en oblique dans la marge gauche,
soulignés, chacun en face de son schéma ; le troisième schéma n'en a
pas. Le $B_G$ de la dernière suite est tel quel sur la page.}

\section{Familles limitées}
\note{« Familles limitées », et au-dessous « cf.\ SGA 6 XIII », sont ses
mots, seuls sur la page 25, qui sert de couverture à ce qui suit.}

\page{26}
Il y a deux variantes : \uncertain{pr}\ill{}

\emph{\uncertain{Assez naturel}} : Si des $F_i$ sont quotients d'un
[$\ill{} = L$], et si \ill{}, les cycles associés, de dimension $d$, de
degrés bornés, alors les $F_i$ sont limités ;

\struck{Les variantes $L = \mathcal{O}(-n)^k$ et \ill{}}

\struck{$L = \ill{}$,} $E = \ill{}$, lorsque \ill{}

$F_i$ \ill{} : \ill{} bornée, \ill{} \ill{} la \ill{} de \struck{\ill{}}
\ill{}, soit $l$. Si $l \le 0$ rien à prouver, \ill{} \ill{} $l$.

\struck{Il faut \ill{} \ill{}} Si $l = 1$, dans les \ill{} \ill{},
\ill{} \struck{$F_i$ \ill{}} \ill{} l'\ill{} \ill{} extension de faisceaux
\ill{} \ill{} limités.

\marginal{(e)} limités, \struck{reste} \ill{}

Si $l = 1$, $F_i \subset \struck{\ill{}}\ \mathcal{O}_{Z_i}$, $Z_i$ un
cycle \uncertain{borné} \ill{} \ill{} de degré \ill{} $L$.
\struck{Le résidu \ill{}} \ill{} $\mathcal{O}_Z$.

\marginal{(d)} \ill{} limités, \ill{} \ill{} fini ; (\ill{} d'un \ill{}
\ill{} faisceaux $L$ dans $\mathcal{O}_Z$).
On \ill{} \ill{} fini, \ill{} \ill{}

\marginal{(f)} Chow \uncertain{général} \ill{}
\note{page d'écriture très rapide, sur sa largeur ; les lettres (d), (e),
(f) cerclées sont dans la marge gauche, en face des lignes où elles sont
données. Seuls les signes mathématiques et quelques mots se lisent.}

\page{27}
faisceaux $F_i$ quotients d'un faisceau donné $E$, et \struck{\ill{}}
à $P_{F_i}$ \uncertain{donnés} \struck{\ill{}} fixés. On \ill{} prouver
que les $(F_i)$ sont limités. \struck{Réciproque}

\struck{\ill{}}
4°) \emph{Résultat de Chow} : les \emph{cycles} de $\mathbf{P}^n_{\mathbf{Z}}$
de degré \struck{\ill{}} borné \uncertain{sont en nombre limité} (i.e.\
\struck{\ill{} fini} \struck{\ill{} fini} \uncertain{fibres} d'une
\struck{\ill{}} $\Gamma \subset \mathbf{P}^n_T \to T$, où $T$ de type fini
sur $\mathbf{Z}$).
\note{« cycles » est doublement souligné. Entre $\mathbf{P}^n$ et la
flèche vers $T$, un indice noirci.}

c) récurrence sur $\sup \dim \operatorname{Supp} F_i = d$, trivial si
$d < 0$. Si $d \ge 0$ on regarde les \ill{}
\[
  0 \to F''_i \to F_i \to F'_i \to 0
\]
où $\dim F''_i < d$ et $F'_i$ \struck{\ill{}} \uncertain{cycles} \ill{}
de dim $d$ \ill{} de $L$, dans un nombre fini \ill{}

\marginal{(d)} Les $F'_i$ sont limités \ill{} $= L''$ \ill{} $= L''$ et
leurs \uncertain{Hilb.} \struck{\ill{} \ill{} dans \ill{}} \ill{}
\uncertain{j'y reviendrai}. Dans les conditions $F_i$ et $F'_i$ limités, les
$P_{F'_i}$ sont en nb fini ; $P_{F''_i} = P_{F_i} - P_{F'_i}$ en nb fini ;
\ill{} les $F''_i$ sont quotients d'un \ill{} et \uncertain{Hilb.} en nb
fini, donc \ill{}

\marginal{(e)} les $F''_i$ sont limités.
\note{(d) et (e) cerclés en marge ; le « c) » qui précède est écrit dans
le corps, non cerclé. La lettre (d) répond à celle de la page 26 : les
deux pages reprennent le même plan.}

\section{[Classes de faisceaux sur les extensions du corps de base]}
\note{titre de l'éditeur, entre crochets.}

\page{28}
\textbf{Définition 1.} Soient $X$ un préschéma \struck{Ag} sur un corps $k$,
$F_i$ \struck{et $G_i$} ($i = 1, 2$) un faisceau algébrique sur
$X_i = X \otimes_k K_i$, où $K_i$ ext.\ de $k$. On dit $F_1$, $F_2$
\emph{$k$-équivalents} s'il existe une extension $K'$ de $k$, des
$k$-homom.\ $u_i : K_i \to K'$, tels que les faisceaux algébriques
$F_i \otimes_{K_i} K'$ sur $X \otimes_k K' \simeq X_i \otimes_{K_i} K'$
soient isomorphes. \struck{La classe d'équivalence} \ill{} de faisceaux.
\note{« $F_i \otimes_{K_i} K'$ » est interligné. Dans la marge gauche, un
petit croquis : une nappe marquée $X$, des verticales, $K'$ ; un trait
vertical borde la définition.}

Les \ill{} \uncertain{géométrique} de \ill{}. On vérifie facilement que
c'est une \uncertain{relation d'équivalence}. Classe d'équivalence de
couples $(K, F)$ \struck{faisceaux} \ill{} : $K$ ext.\ de $k$, $F$
faisceau algébrique sur $X \otimes_k K$ \ill{} pour la transitivité
\ill{} extension \ill{} des $K$ \ill{} \ill{} choisie.

\emph{Si $X$ est de type fini sur $k$,}

\textbf{Proposition 1.} Tout faisceau algébrique sur $X \otimes_k K'$ est
équivalent à un faisceau alg.\ sur $X \otimes_k K$, où $K$ sous-ext.\ de
$K'$ de type fini sur $k$. Si $F_1$ sur $X \otimes_k K_1$ et $F_2$ sur
$X \otimes_k K_2$ sont deux faisceaux équivalents, \struck{\ill{}},
$K_i$ de type fini sur $k$, \struck{de déf.\ \ill{}} \ill{}

\page{29}
dans la définition 1 on peut prendre pour $K'$ une extension de type fini
de $K$.

\emph{Cor.} \ill{} résulte de Chap.\ V, § 1. \emph{Question} :
\uncertain{peut-on prendre} $K'$ \ill{} quelle extension \ill{}
\uncertain{universelle} ?

\emph{Cor.} \ill{} l'ens.\ des classes de faisceaux \ill{}, est en bij.\
\uncertain{canonique} avec \ill{} faisceaux « définis » sur des extensions
de type fini de $k$. [qui est un ens.\ au sens de \ill{}] [\ill{} dans
$L$ = réunion \ill{} $k$, $X$]. \ill{} $\mathfrak{F}(X/k)$.
\note{$\mathfrak{F}(X/k)$ est encadré.}

\struck{\textbf{Définition.} \ill{} corps de base $k$ en $K'$. $E$ \ill{}
$\mathfrak{F}(X/k) \to \ill{}$ \ill{} $X/k$, par $k$.}
\note{bloc barré de trois diagonales ; seuls ces mots s'y lisent.}
\marginal{plus bas \uncertain{complété}}

\ill{} que l'on $\mathfrak{F}(X/k)$ \ill{}, \ill{} \uncertain{foncteur
contravariant} \ill{} si $h$ \ill{} $k$, il \ill{}

\page{30}
\textbf{Corollaire} \uncertain{5}. Soit $\widetilde{K}$ une extension
algébrique\uncertain{ment close} de $k$ \ill{} transcendance \ill{}, et
$\widetilde{X} = X \otimes_k \widetilde{K}$. Alors l'ens.\
$\mathfrak{F}(X/k)$ \struck{iden} s'identifie à un ens.\ quotient de
l'ens.\ des classes de faisceaux \add{coh.}\ $F$ sur $\widetilde{X}$
\ill{} relation : \struck{\ill{}} $F$ et $G$ sont $k$-équivalents ssi il
existe un automorphisme $\sigma$ de $\widetilde{K}/k$ tel que
($\sigma_{\widetilde{X}}$ désignant l'automorphisme
$\mathrm{id}_X \times \sigma$ de $\widetilde{X}$)
$F \simeq \sigma_{\widetilde{X}}(G)$. En d'autres termes
$\mathfrak{F}(X/k)$ s'identifie à l'ens.\ quotient de l'ens.\ des classes
de faisceaux coh.\ sur $\widetilde{X}$ par $\operatorname{Aut}(\widetilde{K}/k)$
y opérant de façon évidente.
\note{« coh. » est interligné.}

Ceci dit, soit $\Omega$ une \ill{}, et $\Phi$ une \ill{} de \ill{}
$\mathfrak{F}(X \otimes_k K)$, \ill{} $X$ \ill{} préschéma algébrique sur
un corps alg.\ clos, \struck{quelconque} \uncertain{fonction}

\marginal{\uncertain{Supposons} $X$ et $Y$ sur $k$ donnés, et pour tt
$K/k$ donné ($K$ alg.\ clos), une partie $\Psi(X \otimes_k K)$ de
$\mathfrak{F}(X \otimes_k K)$ \ill{} extension de base \ill{}}
\note{note écrite en oblique dans l'angle inférieur gauche.}

\page{31}
un ensemble $T(F_1, \ldots, F_n)$ de \uncertain{classes}
\uncertain{choisies} de faisceaux (\ill{}) sur $X$, de façon que :

(i) si les $F_i$ sont remplacés par des faisceaux $F'_i$ isomorphes,
$T(F'_1, \ldots, F'_n)$ \uncertain{est inchangé} \struck{: $T(E\ill{})$}
un faisceau $G \in T(F_1, \ldots, F_n)$ ;

(ii) la compatibilité avec l'extension du corps de base :
$T(F_1 \otimes k', \ldots, F_n \otimes k')$ \uncertain{est} (à isomorphie
près) l'ens.\ des faisceaux $G \otimes k'$ où $G \in T(F_1, \ldots, F_n)$.
\marginal{(i) \uncertain{Généralisation}}

\struck{\textbf{Déf.}} Si alors $A_i$ sont des parties
(\uncertain{$i = 1, \ldots, n$}) de $\mathfrak{F}(X/k)$, on désigne par
$\mathfrak{F}(A_1, \ldots, A_n)'$ \uncertain{l'ens.}\ des classes
\uncertain{dans} $\mathfrak{F}(X/k)$ des cl.\ des $T(F_1, \ldots, F_n)$,
où $F$ \uncertain{sont} \ill{} sur $\widetilde{X}$ et
$\mathrm{cl}_K(F_i) \in A_i$ pour tt $i$.
\note{« Déf. » est encadré et biffé.}

On voit \struck{\ill{}} facilement que la définition ne dépend pas du
choix de $\widetilde{K}$ ; en effet si \ill{} une extension
$\widetilde{K}'$ de $\widetilde{K}$ \ill{} ayant les mêmes propriétés, et
il en \ill{} facile de démontrer \ill{}. D'ailleurs, facile de \ill{}
direct de $\mathfrak{F}(A_1, \ldots, A_n)$ \ill{}, [définition \ill{}
$\widetilde{K}$, \ill{} l'un des \ill{}]
\marginal{\uncertain{propriété}}
\note{la page, d'écriture rapide, prolonge la page 30 : le $T$ est
l'opération « de $n$ faisceaux » annoncée au bas de celle-ci. Une note
oblique dans la marge inférieure gauche ne se lit pas.}

\page{32}
$\mathfrak{F}(F_1, \ldots, F_n)$, où les $F_i$ sont des faisceaux sur
$\widetilde{X} = X \otimes_k K$, \add{$K$ alg.\ clos,} $\mathrm{cl}_k F_i
\in A_i$.
\note{« $K$ alg. clos » est interligné.}

D'ailleurs, si on prend \struck{\ill{}} \ill{} pour définir des faisceaux
$F_i$ sur des extensions \emph{distinctes} $K_i/k$, $F_i$ de classe
$\xi_i \in \mathfrak{F}(X/k)$, on obtient $\mathfrak{F}(\xi_1, \ldots,
\xi_n)$ en prenant de \ill{} \ill{} qui \uncertain{fait} \ill{} une
extension composée des $K_i$ sur $k$, et désignant \struck{$K$} $L_\alpha$,
et \ill{} l'une des \ill{} en \ill{} de \ill{} de $\mathfrak{F}(X/k)$,
\struck{$\mathfrak{F}(F_1^{(\alpha)}, \ldots, F_n^{(\alpha)})$} \ill{}
$\mathfrak{F}(F_1 \otimes_k L_\alpha, \ldots, F_n \otimes_k L_\alpha)$.

Si $\mathfrak{F}$ est défini aussi pour des faisceaux sur $X_K$,
\uncertain{seulement} $K$ alg.\ clos, \ill{} satisfaisant
\uncertain{correspondantes} des (i) et (ii), \ill{} \ill{} \ill{} \ill{}
les \struck{\ill{}} \ill{} complications précédentes de \ill{} $K$ alg.\
clos, et \struck{prenant} \ill{} ; \ill{} obtenons \ill{} algébriquement un
\uncertain{extension}, \ill{} \uncertain{correspondants}.

\marginal{pour 2 faisceaux $F_i$ \uncertain{remplacés} \ill{} \ill{}
\ill{} \ill{} ! \ill{} $X_i$ \ill{} \ill{}}
\note{note oblique dans la marge gauche, presque entièrement illisible.}

\page{33}
\emph{Exemples.}

a) foncteurs $\mathcal{H}om$, $\otimes$, et leurs dérivés.
[2 arguments, $X = Y$] et leurs \uncertain{conjugués} en \ill{}

(d) Si $F$ et $G$ sont \ill{} des faisceaux, \ill{} l'ensemble
$N(F, G)$ des \ill{} \ill{} [\ill{} $=$ \ill{} \ill{}, kif kif, \ill{}
\ill{} \ill{} identifié \ill{} précédemment \ill{}]. [2 arg., $X = Y$]

(b) F\ill{} \ill{} \ill{} \ill{} (1 arg., $f : X \to Y$ donné).

(c) F\ill{} \ill{} \ill{} et foncteurs dérivés (1 arg., $f : X \to Y$
\struck{\ill{}} donné).

(e) Ens.\ des extensions de $F$ par $G$ (2 arg., $X = Y$).

(f) Ens.\ des \ill{} \ill{} : 2 \ill{} \ill{} \ill{}, \ill{} ;
décroissante, de $F$ (1 arg., $X = Y$).

(f) \add{$F$.\ill{}} \struck{Ens.\ des \ill{}}
$\mathcal{O}_{(\mathrm{S}\ill{}\, F)_{\mathrm{red}}}$ ($K$ alg.\ clos, 1
arg., $X = Y$).

(h) Ens.\ des \struck{\ill{}} \struck{faisceaux} composantes primaires de
$F$ associées aux composantes de son support ($K$ alg.\ clos, \ill{},
$X = Y$).

\struck{(i) \ill{}}
\note{les lettres sont les siennes, dans cet ordre : a, (d), (b), (c),
(e), (f), (f), (h) ; (d), (e) et le second (f) sont précédés d'un
astérisque, et des traits dans la marge gauche relient (d) à (e) et (b)
à (f). La dernière ligne, (i), est entièrement biffée.}

\section{[Stabilité des familles limitées ; Tor et changement de base]}
\note{titre de l'éditeur, entre crochets.}

\page{34}
\marginal{(a)}
Soient $A, B \subset \mathfrak{F}(X/S)$, $A$ et $B$ limités, alors
$A \otimes B$ et $\mathcal{H}om(A, B)$ sont limités, \ill{} \ill{} \ill{}
$\mathcal{T}or_i(A, B)$ et $\mathcal{E}xt^i(A, B)$.

\struck{Dém.} \ill{}, supposons
\[
  \begin{cases}
    A = \struck{\mathfrak{F}}\ \mathcal{E}(F), & F \text{ \uncertain{défini} sur } X \times_S S' \\
    B = \mathcal{E}(G), & G \text{ --- } X \times_S S''
  \end{cases}
\]

\struck{\ill{} \ill{}} \emph{Avec} \uncertain{ces notations},
$H = F \otimes_X G$ sur $X \times_S S'''$, $S''' = S' \times_S S''$.
Alors tt \uncertain{type} d'extension \uncertain{composée} d'un $k(s')$ et
$k(s'')$, sur $k(s)$, [$s' \mapsto s$, $s'' \mapsto s$ ;
$s' \in S'$, $s'' \in S''$] provient d'un (\uncertain{unique})
$s''' \in S'''$, $s''' \mapsto s'$, $s''' \mapsto s''$ ;
\ill{} \ill{} \ill{}
\[
  (F_{s'} \otimes_{k(s')} K) \otimes_{\struck{X}\, X_{s'} \otimes_{k(s)} K}
  (G_{s''} \otimes_{k(s'')} K) \simeq H_{s'''} .
\]
\note{dans la marge gauche, les huit sommets d'un cube, $X$, $X'$, $X''$,
$X'''$ au-dessus de $S$, $S'$, $S''$, $S'''$, reliés par des traits sans
flèche.}

Comme $S'''$ est de type fini sur $S$, la \uncertain{conclusion} \ill{}
$\otimes$ s'ensuit, on a bien prouvé
\[
  \struck{\ill{}\ \mathcal{E}(F) \otimes}\ \mathcal{E}_{X/S}(F) \otimes
  \mathcal{E}_{X/S}(G) \subset \mathcal{E}_{X/S}(H) .
\]

Pour $\mathcal{H}om$, \ill{} \struck{\ill{} \uncertain{supposons} $X'$ et
$F$ $S'$-plats} \ill{} ; $\bar{F} = F \otimes_{S'} S'''$,
$\bar{G} = G \otimes_{S'} S'''$ et \ill{} \ill{} [\ill{}] \ill{} on
\uncertain{trouve} \ill{} une partition \ill{} de $S'''$ en \ill{}
$\to T_i$ \ill{} ; posant $F_i = \bar{F} \otimes_{S'''} T_i =
F \otimes_{S'} T_i$, $G_i = \bar{G} \otimes_{S'''} T_i =
G \otimes_{S''} T_i$, $X_i^{*} = X''' \times_{S'''} T_i = X \times_S T_i$,
\note{dans $\bar{G} = G \otimes_{S'} S'''$, l'indice $S'$ est tel quel sur
la page, là où l'on attend $S''$ (comme dans $G \otimes_{S''} T_i$
ensuite).}

\page{35}
la formation de $\mathcal{H}om_{X_i}(F_i, G_i)$ commute à
\uncertain{tt} changement de base $T' \to T_i$ \ill{} \ill{} extension de
base [\uncertain{pour ceci}, écrivons \ill{} \ill{}
\uncertain{présentation} libre de type fini
$L_1 \to L_0 \to \bar{F} \to 0$,
\struck{\ill{} \ill{} $0 \to Z_1 \to L_1 \to Z_0 \to 0$}
\[
  0 \to Z_0 \to L_0 \to \bar{F} \to 0
\]
et \struck{\ill{}} \ill{}, $Z_i$ \uncertain{localement} \uncertain{libres}
\ill{} et \ill{} \ill{} \ill{}
$\operatorname{Ker}\bigl(\mathcal{H}om(L_0, F) \to
\mathcal{H}om(L_1, F)\bigr)$ \ill{} \ill{}.
\note{la suite $0 \to Z_0 \to L_0 \to \bar{F} \to 0$ est encadrée.}

\ill{} \ill{} \ill{} \ill{} \uncertain{convenable} \ill{} \ill{} \ill{}
de $S'''$ \ill{} \ill{} $T_i$, dans \ill{} \ill{} \ill{} \ill{}
\uncertain{décomposition} : $T_i$, la formation \ill{}, je \ill{} \ill{}
\ill{} \ill{} \ill{} \ill{} Ker \ill{} \ill{} \ill{} \ill{} base \ill{}
\uncertain{composé}. \ill{} \ill{} \ill{} \ill{} \ill{}.]

Le cas des $\mathcal{T}or_i$ et $\mathcal{E}xt^i$ \ill{} \ill{} \ill{},
\ill{} \ill{} \ill{} \ill{} \ill{}.

\emph{Lemme.} Soit $Y$ de type fini sur $T$ \uncertain{noethérien},
$K^{\bullet}$ un \ill{} \ill{} \ill{} faisceaux cohérents, \ill{} \ill{}
\ill{} de degré $+1$ \ill{} \ill{} \ill{} \ill{} \ill{}, $N$ \ill{}
\ill{}. Mais il y a une décomposition de $T$ en $T_i$ loc.\ fermés tels
que
\marginal{\ill{} $L_i$ \ill{} \uncertain{localement} \ill{} \ill{} IV,
§ 3}
\note{note oblique dans la marge gauche ; seul « IV, § 3 » se lit
sûrement.}

\page{36}
\struck{\ill{}} \ill{} \uncertain{posant}
$K^{(i)} = K \times_T T_i$, \uncertain{pour} \struck{\ill{}}
$T' \to T_i$, \ill{} \ill{} \ill{} \uncertain{homomorphismes} canoniques
\[
  H^q(K^{(i)} \otimes_{T_i} T') \longleftarrow H^q(K^{(i)}) \otimes_{T_i} T'
\]
\uncertain{soient} \ill{} \uncertain{isomorphismes} pour $|q| \le N$.

Démonstration, \ill{} \ill{} \ill{} \ill{} \ill{}
\uncertain{nilpotent}\ill{} \ill{} $T$, \ill{} \ill{} \ill{} \ill{} :
\ill{} \ill{} \ill{} \ill{} \ill{} \ill{} \ill{} $U$ \ill{} \ill{} \ill{}
$T$, $T$ \uncertain{est} \uncertain{réduit} \ill{} [\uncertain{ouvert}
\ill{} \ill{} \ill{} et $T_i$, (\ill{} \ill{}, \ill{}] \ill{} \ill{} \ill{}
\ill{} [\uncertain{choisi} \ill{}] \uncertain{tels} \ill{} \ill{} \ill{}
$B_i \otimes_T U$ \ill{}
\[
  K^{i-1} \xrightarrow{\ d^{i-1}\ } K^i \to K^{i+2}
\]
\struck{$(\operatorname{Coker} d^i) \times_T U$ \ill{} \ill{} \ill{}}
\ill{} si $i \le N+1$ \struck{[$U$ \ill{} \ill{} \ill{} \ill{} \ill{}]}
$U$\ill{} $U$-\ill{}, pour $|i| \le N+1$ \ill{}
\note{la flèche de droite va bien de $K^i$ à $K^{i+2}$ sur la page.}

\marginal{(b)}
Soit $f : X \to Y$ un $S$-morphisme de préschémas de type finis. Soit
$A \subset \mathfrak{F}(Y/S)$ un \uncertain{ensemble} limité, alors
$f^{*}(A)$ et les $\mathcal{T}or_{*}^{\mathcal{O}_Y}(A, \mathcal{O}_X)$
sont limités.

Il y a \uncertain{un énoncé} analogue \ill{} : \ill{} \ill{} \ill{}
proposition \ill{} \ill{} le produit, \ill{} $X$ et $Z$ sur $Y$,
\struck{\ill{}} $A \in \mathfrak{F}(X/S)$, $B \in \mathfrak{F}(Z/S)$, et
les $\mathcal{T}or_i^{\mathcal{O}_Z}(A, B)$.
\note{l'indice de ce dernier $\mathcal{T}or$ se lit $\mathcal{O}_Z$ ;
on attendrait $\mathcal{O}_Y$. Laissé tel quel.}

\page{37}
Soient $X_i$ \uncertain{sur} $Y$ ($1 \le i \le n$), des $F_i$ sur les
$X_i$, \struck{\ill{}} un \uncertain{changement} de base $Y' \to Y$,
\struck{\ill{}} \ill{} \uncertain{homomorphisme}
\[
  \mathcal{T}or_{*}^{\mathcal{O}_Y}(F_1, \ldots, F_n) \otimes_Y Y'
  \quad\text{et}\quad
  \mathcal{T}or_{*}^{\mathcal{O}_{Y'}}(F'_1, \ldots, F'_n)
\]
\note{dans la marge gauche, en colonnes, les lettres $X_1\ X_2$,
$X'_1\ X'_2$, $Y$, $Y'$, $S$, $S'$, sans traits.}

On \uncertain{sait} \ill{} \ill{} \ill{} \ill{} un homomorphisme canonique
du \uncertain{premier} dans le second. À quelle condition \ill{} \ill{}
\ill{} un isomorphisme ?

C\struck{\ill{} \ill{} \ill{} \ill{} \ill{} \ill{} \ill{}} \ill{},
\uncertain{limités} \uncertain{pour} tous (\ill{} \ill{} \ill{} \ill{})
\uncertain{que}
\[
  \mathrm{Tor}_p^{\mathcal{O}_S}\bigl(\mathrm{Tor}_q^{\mathcal{O}_Y}(F_1, \ldots, F_n)
  \otimes_{\ill{}} \ldots, \mathcal{O}_{S'}\bigr)
\]
il suffit que les $\mathrm{Tor}_q^{\mathcal{O}_Y}(F_1, \ldots, F_n)$
soient $Y$-plats.
\note{une ligne entière barrée ; au-dessus, interligné, « de = \ill{} ».
Le signe entre les deux $\mathrm{Tor}$ est surchargé d'une tache.}

\page{38}
\emph{Proposition.} $X_i$ ($1 \le i \le n$) de type fini sur $Y$ de type
fini sur $S$ noeth., $F_i$ coh.\ sur $X_i$, $N$ un entier. Alors il existe
un \struck{\ill{}} \uncertain{ouvert dense} $U$ de $S$, tel que,
\uncertain{posant} $S' = U_{\mathrm{red}}$, $X'_i = X_i \times_S S'$,
$Y' = Y \times_S S'$, $F'_i = F_i \otimes_S S'$, on ait ceci : pour tout
changement de base $T \to S'$, les homomorphismes canoniques
\[
  \mathrm{Tor}_p^{\mathcal{O}_{S'}}(F'_1, \ldots, F'_n) \otimes_{S'} T
  \longrightarrow
  \mathrm{Tor}_p^{\mathcal{O}_T}(F'_{1T}, \ldots, F'_{nT})
\]
(où $F'_{\alpha T} = F'_\alpha \otimes_{S'} T$) soient des
\emph{isomorphismes} pour $p \le N$.
\note{un trait vertical borde l'énoncé. Dans $S' = U_{\mathrm{red}}$, le
$S$ et les $S$ des produits fibrés sont récrits sur d'autres lettres
noircies. Les $\mathrm{Tor}$ sont pris sur $\mathcal{O}_{S'}$ et
$\mathcal{O}_T$ tels qu'écrits.}

On utilise \ill{} \uncertain{ceci} \ill{}
\add{Avec les notations précédentes mais \ill{} conditions de finitude,}
\struck{\ill{} \ill{} \ill{} \ill{}}
\note{l'ajout est interligné et barré avec la ligne qu'il surmonte ; la
page 41 (lot 3) reprend « Avec les notations précédentes, mais sans
conditions de finitude ».}

\marginal{Faut-il \ill{} ? \ill{}}
\emph{Lemme} (\uncertain{Important}\ill{}). Supposons que $Y$ soit plat
sur $S$, \uncertain{et} les $X_i$ plats sur $S$, et que les
$\mathrm{Tor}_i^{\mathcal{O}_Y}(F_1, \ldots, F_n)$ pour $i \le N$ soient
plats sur $S$. Alors pour \ill{} changement de base $T \to S$,
l'homomorphisme canonique
\[
  \mathrm{Tor}_i^{\mathcal{O}_S}(F_1, \ldots, F_n) \times_S T \longrightarrow
  \mathrm{Tor}_i^{\mathcal{O}_T}(F_{1(T)}, \ldots, F_{n(T)})
\]
est un isomorphisme pour $i \le N$.
\note{au-dessus de « Supposons », un ajout interligné illisible. Le
premier $\mathrm{Tor}$ du lemme porte $\mathcal{O}_Y$, celui de la flèche
$\mathcal{O}_S$ : tels qu'écrits.}

\page{39}
\struck{\ill{}} Pour la \uncertain{preuve}, on \uncertain{peut}
supposer \ill{} $S$, $Y$, les $X_i$, \uncertain{affines} d'anneaux $A$,
$B$, $B_i$, $A'$, ($F_i$ \uncertain{définis} par \struck{$M$} des
$B_i$-modules $M_i$. Soit $L^{(i)}$ une résolution libre de $M_i$
\uncertain{par} $B_i$-modules \ill{}, \uncertain{posons}
\[
  K = \bigotimes_{\substack{B \\ 1 \le i \le n}} L^{(i)} .
\]

Alors \struck{\ill{}}, les $L^{(i)}$ sont \emph{\uncertain{plats} sur $A$},
$M_i$ \uncertain{est} \struck{plat} sur $A$, il en \uncertain{résulte} que
$L'^{(i)} = L^{(i)} \otimes_A A'$ est une résolution de
$M'_i = M_i \otimes_A A'$, \uncertain{considéré} comme $B'_i$-module (où
$B'_i = B_i \otimes_A A'$). D'où la \uncertain{conclusion} que
\[
  K' = K \otimes_A A' = \bigotimes_{\substack{B' \\ 1 \le i \le n}}
  L^{(i)} \otimes_B B' = \bigotimes_{\substack{B' \\ 1 \le i \le n}} L^{(i)\prime}
\]
\ill{} \ill{} que $\mathrm{Tor}_{*}^{B'}(M'_1, \ldots, M'_n)$. Or
$K^{\bullet}$ \uncertain{est un complexe} $A$-plat,
\uncertain{on a donc} une suite spectrale
\[
  H_{*}(K \otimes_A A') \Longleftarrow \mathrm{Tor}_p^A(H_q(K), A')
\]
i.e.
\[
  \mathrm{Tor}_{*}^{B'}(M'_1, \ldots, M'_n) \Longleftarrow
  \mathrm{Tor}_p^A\bigl(\mathrm{Tor}_q^B(M_1, \ldots, M_n), A'\bigr) .
\]
Il \uncertain{suffit} donc que les $\mathrm{Tor}_q^{\ill{}}(M_1, \ldots,
M_n)$ \uncertain{soient} $A$-plats pour $q \le N$, \uncertain{pour}
l'\uncertain{homomorphisme} \ill{} \ill{} \ill{} \ill{}
\note{dans la chaîne $K' = \ldots$, un terme intermédiaire est noirci. Dans
la marge gauche, en colonnes, $M_1\ M_2$, $M'_1\ M'_2$ ; $B_1\ B_2$,
$B'_1\ B'_2$ ; $B$, $B'$ ; $A$, $A'$, cernés d'un trait courbe ; plus
bas, une note oblique : « \ill{} \ill{} \ill{} spectrale, i.e.\ \ill{}
\ill{} \ill{} ».}

\page{40}
\[
  \mathrm{Tor}_p^B(M_1, \ldots, M_n) \otimes_A A' \longrightarrow
  \mathrm{Tor}_p^{B'}(M'_1, \ldots, M'_n)
\]
est un isomorphisme pour $p \le N$. On \ill{} \uncertain{ainsi} \ill{}
la proposition.

\struck{Si \ill{} les conditions de la \uncertain{prop.}}

\textbf{Corollaire.} Il existe une décomposition de $S$ en \ill{} loc.\
fermés $S_\alpha$ \uncertain{en nombre fini}, ($1 \le i \le n$) tels que,
posant $X_i^{(\alpha)} = X_i \times_S S_\alpha$,
\add{$Y^{(\alpha)} = Y \times_S S_\alpha$,}
$F_i^{(\alpha)} = F_i \otimes_S S_\alpha$, \struck{on ait \ill{}} \ill{}
pour \ill{} $T \to S_\alpha$, l'homomorphisme canonique
\[
  \mathrm{Tor}_p^{\mathcal{O}_{S_\alpha}}(F_1^{(\alpha)}, \ldots, F_n^{(\alpha)})
  \otimes_{S_\alpha} T \longrightarrow
  \mathrm{Tor}_p^{\mathcal{O}_T}(F^{(\alpha)}_{1(T)}, \ldots, F^{(\alpha)}_{n(T)})
\]
soit un isomorphisme pour $|p| \le N$.

\textbf{Proposition.} $X$ de type fini sur $S$ noeth., $F$ et $G$
\add{coh.}\ faisceaux \uncertain{cohérents} sur $X$. \add{$N$ un entier.}
Il existe \ill{} un \uncertain{ouvert dense} $U$ de $S$, telle que,
posant $S' = U_{\mathrm{red}}$, $X' = X \times_S S'$,
$F' = F \times_S S'$, $G' = G \times_S S'$, on ait ceci : pour tout
changement de base $T \to S$, l'homomorphisme canonique
\[
  \mathcal{E}xt^p_{\mathcal{O}_{S'}}(F', G') \otimes_{S'} T \longrightarrow
  \mathcal{E}xt^p_{\mathcal{O}_T}(F'_{(T)}, G'_{(T)})
\]
soit un isomorphisme pour $|p| \le N$.
\note{« $N$ un entier » est interligné. Les deux énoncés sont bordés d'un
trait vertical. La page 41 (lot 3) continue avec le lemme correspondant
pour les $\mathcal{E}xt$.}

\end{document}
